\SourceLeaf{062}
\begin{center}
{\large SUPPLEMENTUM\par}
\medskip THEORIAE\par
\medskip {\large COMBINATIONIS OBSERVATIONUM\par}
\medskip ERRORIBUS MINIMIS OBNOXIAE.\par
\medskip\rule{.25\linewidth}{.3pt}
\end{center}
\Article{1}
In tractatione theoriae combinationis observationum Volumini V Commentationum Recentiorum inserta supposuimus, quantitates eas, quarum valores per observationes praecisione absoluta non gaudentes propositi sunt, a certis elementis incognitis ita pendere, ut in forma functionum datarum horum elementorum exhibitae sint, reique cardinem in eo verti, ut haec elementa quam exactissime ex observationibus deriventur.

In plerisque quidem casibus suppositio ista immediate locum habet. In aliis vero casibus problematis conditio paullo aliter se offert, ita ut primo aspectu dubium videatur, quonam pacto ad formam requisitam reduci possit. Haud raro scilicet accidit, ut quantitates eae, ad quas referuntur observationes, nondum exhibitae sint in forma functionum certorum elementorum, neque etiam ad talem formam reducibiles videantur, saltem non commode vel sine ambagibus: dum, ex altera parte, rei indoles quasdam conditiones suppeditat, quibus valores veri quantitatum observatarum exacte et necessario satisfacere debent.

Attamen, re propius considerata, facile perspicitur, hunc casum ab altero revera essentialiter haud differre, sed ad eundem reduci posse. Designando scilicet multitudinem quantitatum observatarum per \(\pi\), multitudinem aequationum conditionalium autem per \(\sigma\), eligendoque e prioribus \(\pi-\sigma\) ad lubitum, nihil impedit, quominus has ipsas pro elementis accipiamus, reliquasque, quarum mul\SourceLeaf{063}%
titudo erit \(\sigma\), adiumento aequationum conditionalium tamquam functiones illarum consideremus, quo pacto res ad suppositionem nostram reducta erit.

Verum enim vero etiamsi haec via in permultis casibus satis commode ad finem propositum perducat, tamen negari non potest, eam minus genuinam, operaeque adeo pretium esse, problema in ista altera forma seorsim tractare, tantoque magis. quod solutionem perelegantem admittit. Quin adeo, quum haec solutio nova ad calculos expeditiores perducat, quam solutio problematis in statu priore, quoties \(\sigma\) est minor quam \(\frac12\pi\), sive quod idem est, quoties multitudo elementorum in commentatione priore per \(\rho\) denotata maior est quam \(\frac12\pi\), solutionem novam, quam in commentatione praesenti explicabimus, in tali casu praeferre conveniet priori, siquidem aequationes conditionales e problematis indole absque ambagibus depromere licet.

\Article{2}
Designemus per \(v,v',v''\) etc. quantitates, multitudine \(\pi\), quarum valores per observationem innotescunt, pendeatque quantitas incognita ab illis tali modo, ut per functionem datam illarum, puta \(u\), exhibeatur: sint porro \(l,l',l''\) etc. valores quotientium differentialium
\[
\frac{\mathrm du}{\mathrm dv},\quad \frac{\mathrm du}{\mathrm dv'},\quad \frac{\mathrm du}{\mathrm dv''}\ \text{etc.}
\]
valoribus veris quantitatum \(v,v',v''\) etc. respondentes. Quemadmodum igitur per substitutionem horum valorum verorum in functione \(u\) huius valor verus prodit, ita, si pro \(v.\ v',v''\) etc. valores erroribus \(e,e',e''\) etc resp. a veris discrepantes substituuntur, obtinebitur valor erroneus incognitae, cuius error statui potest
\[
=le+l'e'+l''e''+\text{etc.}
\]
siquidem, quod semper supponemus, errores \(e,e',e''\) etc. tam exigui sunt, ut (pro functione \(u\) non lineari) quadrata et producta negligere liceat. Et quamquam magnitudo errorum \(e,e',e''\) etc. incerta maneat, tamen incertitudinem tali incognitae determinationi inhaerentem generaliter aestimare licet, et quidem per errorem medium in tali determinatione metuendum, qui per principia commentationis prioris fit
\[
=\BareRoot{(llmm+l'l'm'm'+l''l''m''m''+\text{etc.})}
\]
\SourceLeaf{064}%
denotantibus \(m,m',m''\) etc. errores medios observationum, aut si singulae observationes aequali incertitudini obnoxiae sunt,
\[
=m\BareRoot{(ll+l'l'+l''l''+\text{etc.})}
\]
Manifesto in hoc calculo pro \(l,l',l''\) etc. aequali iure etiam eos valores quotientium differentialium adoptare licebit, qui valoribus observatis quantitatum \(v,v',v''\) etc. respondent.

\Article{3}
Quoties quantitates \(v,v',v''\) etc. penitus inter se sunt independentes, incognita unico tantum modo per illas determinari poterit: quamobrem tunc illam incertitudinem nullo modo nec evitare neque diminuere licet, et circa valorem incognitae ex observationibus deducendum nihil arbitrio relinquitur.

At longe secus se habet res, quoties inter quantitates \(v,v',v''\) etc. mutua dependentia intercedit, quam per \(\sigma\) aequationes conditionales
\[
X=0,\quad Y=0,\quad Z=0\ \text{etc.}
\]
exprimi supponemus, denotantibus \(X,Y,Z\) etc. functiones datas indeterminatarum \(v,v',v''\) etc. In hoc casu incognitam nostram infinitis modis diversis per combinationes quantitatum \(v,v',v''\) etc. determinare licet, quum manifesto loco functionis \(u\) adoptari possit quaecunque alia \(U\) ita comparata, ut \(U-u\) indefinite evanescat, statuendo \(X=0,\ Y=0,\ Z=0\) etc.

In applicatione ad casum determinatum nulla quidem hinc prodiret differentia respectu valoris incognitae, si observationes absoluta praecisione gauderent: sed quatenus hae erroribus obnoxiae manent, manifesto in genere alia combinatio alium valorem incognitae afferet. Puta, loco erroris
\[
le+l'e'+l''e''+\text{etc.}
\]
quem functio \(u\) commiserat, iam habebimus
\[
Le+L'e'+L''e''+\text{etc.}
\]
si functionem \(U\) adoptamus, atque valores quotientium differentialium \(\frac{\mathrm dU}{\mathrm dv},\frac{\mathrm dU}{\mathrm dv'},\frac{\mathrm dU}{\mathrm dv''}\) etc. resp. per \(L,L',L''\) etc. denotamus. Et quamquam errores ipsos assignare nequeamus, tamen errores medios in diversis observationum combinationibus me%
